\subsection{Structural reductions at fifteen arrows}

\begin{lemma}\label{C:connected-tree-budget}
A connected finite category on $k$ objects has at least $2k-1$
arrows.  If equality holds, its classifying space is contractible.
Consequently a fifteen-arrow category realizing $\mathbb{CP}^2$
has at most seven objects.
\end{lemma}
\begin{proof}
There are $k$ identity arrows, and the connected underlying
undirected graph needs $k-1$ cross arrows on a spanning tree.
In the equality case these are all arrows: every tree edge has
one orientation, there are no other endomorphisms, and there
are no parallel or reverse cross arrows.  A composable pair
of nonidentity arrows would force a third edge on a cycle of
this tree.  Thus no such pair exists.  The ordinary nerve is
the tree itself and is contractible.  For $k\geq8$ the stated
fifteen-arrow consequence follows from the inequality and its
equality case.
\end{proof}

\begin{lemma}\label{C:nine-bounded-rank-one}
Let $|M|\leq9$, let $S$ be a nonempty finite right $M$-set,
and suppose the act cone $X=B\mathcal C(M,S)$ has bounded rational
homology and $\dim H^2(X;\mathbb Q)\leq1$.  Every degree-two square
is zero.  No assumption on $H^1(X;\mathbb Q)$ is needed.
\end{lemma}
\begin{proof}
This records exactly the obstructions in the proof of
Theorem~\ref{C:nine-act-cones}, without its unnecessary target
specialization.  Through order eight use
Lemma~\ref{C:eight-bounded-cups}.  At order nine, zero vertex
$H^2$ makes every degree-two class relative and its square zero.
Two minimal rows give Theorem~\ref{C:two-row-act-cups}.
The projective-dimension-two cases have $H_4(X;\mathbb Q)=0$.
The exact periodic and stable-retract cases make nonzero fourth
homology recur in arbitrarily high degrees, contradicting the
stated bound.  The seven remaining row modules split into
characters, so Proposition~\ref{C:split-row-act-cups} excludes a
nonzero square when $H^2$ has dimension at most one.
These are all rows of the independently verified order-nine
partition.  None of these arguments uses vanishing first homology.
\end{proof}

\begin{proposition}\label{C:fifteen-three-SCC}
A three-object category with at most fifteen arrows and a
nontrivial strongly connected component does not realize
$\mathbb{CP}^2$.
\end{proposition}
\begin{proof}
Deleting an object which is a retract of another removes its
identity and at least two opposite cross arrows, hence leaves
at most twelve arrows.  The proved smaller exclusion applies.
Assume therefore that no such retract exists.

If all three objects form a strongly connected component,
the categorical Rees description gives finite-group homotopy
type when all row counts or all column counts are one.
Otherwise both total counts are at least four, so the minimal
ideal alone has at least sixteen arrows.  Both alternatives
exclude the proposed type.

It remains to consider a strong two-object block $S$, a singleton
monoid $M$, and cross arrows in one direction.  Write $s,m,t$
for their sizes.  Then $s\geq6$, $t\geq2$, and $s+m+t\leq15$.
If $t=2$, there is one cross arrow at each object of $S$.
The category is the categorical join of $S$ with $M$ (in one
of the two orders).  Its classifying space is a join and has
zero positive-degree rational products.
We may suppose $t\geq3$.

If $BS$ is not of finite-group type, its minimal ideal has at
least nine arrows and the two identities lie outside it.
Thus $s\geq11$, forcing $(s,m,t)=(11,1,3)$.
The minimal group is trivial, the total row and column counts
are three, and only the two identities lie outside the ideal.
Proposition~\ref{C:eleven-strong-three-exit}, or its opposite,
excludes this case.

Otherwise $BS$ is rationally acyclic.  If $m\leq4$, both
vertex-block spaces in the correspondence-bar pushout are
rationally acyclic.  Collapsing them to points gives a rational
cohomology-ring equivalence to a suspension, whose products vanish.
If $m\geq5$, the inequalities $t\geq3$ and $s+m+t\leq15$
give $s\leq7$ and $m\leq6$.  As in
Lemma~\ref{C:fourteen-three-SCC}, the minimal Rees data of $S$
have trivial group and one row and one column at each object.
Both constant rational $S$-modules are projective.
The actual orbit-collapse functor therefore identifies the
whole category rationally, including its cohomology ring,
with a finite act cone over $M$.  Its rational homology is
bounded and its second cohomology is one-dimensional if the
original category has the proposed type.  Apply
Lemma~\ref{C:nine-bounded-rank-one}.
\end{proof}

\begin{lemma}[Star reductions at fifteen arrows]
\label{C:fifteen-star-deletions}
For an outward three-object star of total size fifteen, let
$(m,a,b,x,y)$ denote the central and leaf endomorphism orders
and the two leg sizes.  The generic star exclusions in
Lemma~\ref{C:generic-star-deletions} still hold, with each
bound eight there replaced by nine.  There is also an exclusion
when $m\leq4$ and $x,y\leq2$.  The same assertions hold for
inward stars.  If every vertex has zero rational $H^2$, the
star is excluded as before.
\end{lemma}
\begin{proof}
The singleton-leg, join, and relative-product arguments in the
earlier lemma do not depend on its total-size bound.  For a trivial
center the nerve is a wedge of the two proper collages; each has
at most thirteen arrows, because the other leaf and leg each
contribute at least one arrow.  A wedge of connected spaces of
type $\mathbb{CP}^2$ must have a projective-plane summand: the
free-product formula for $\pi_1$, the direct sum in reduced
homology, and the primitive nonzero degree-two square show
that the other summand is simply connected and acyclic, hence
contractible.  The smaller-category exclusion rules this out.

For $m\leq9$ and $a,b\leq2$, exact rational coinvariants at
the two leaves replace them by trivial monoids and replace the
legs by two actual nonempty finite central acts $S_1,S_2$.
The quotient functor is a rational homology equivalence on each
vertex and edge bar, hence an isomorphism of rational cohomology
rings on the star.  Identify the two sink vertices of this new
star.  On nerves this gives precisely the act cone for
$S_1\sqcup S_2$, since no nonidentity arrow exits a sink.
On realizations it identifies two points in a connected CW space;
the resulting space is homotopy equivalent to the star wedged
with $S^1$.  Thus a hypothetical projective-plane star would
give an act cone with bounded rational homology, one-dimensional
$H^2$, and a nonzero square.  This contradicts
Lemma~\ref{C:nine-bounded-rank-one}.

For the earlier distinguished-leaf case with $m\leq2$, other
leaf order at most four, and distinguished leaf order at most
nine, collapse the rationally acyclic central vertex.
This is a rational cohomology-ring equivalence, and its quotient
is the wedge of the two collage cofibers.  The other collage
has two rationally acyclic vertices and zero products.
Exact central coinvariants turn the distinguished collage
rationally into a finite act cone over its leaf monoid.
Bounded rational homology and dimension at most one for its
$H^2$ follow from the wedge decomposition of a hypothetical
projective-plane star.  The same lemma excludes its square.

Finally, if $m\leq4$ and both legs have at most two arrows,
collapse the rationally acyclic center in the same way.
Each of the two collage summands has a rationally acyclic
vertex and at most two cross arrows, so
Lemma~\ref{C:two-object-cup-pruning} gives zero positive products
on each summand.  Their wedge has zero products as well.
Opposition proves the inward assertions.
\end{proof}
